\section{Data and Variables}
\label{sec:data}

\subsection{Data Sources}

This study assembles a panel of 338 prefecture-level cities over
2005--2023, drawing on four categories of data.

\paragraph{MODIS NDVI.}
We retrieve 16-day MODIS MOD13Q1 NDVI composites from the Microsoft
Planetary Computer STAC API at 250\,m spatial resolution. For each
city, we compute the spatial mean NDVI within a 15\,km buffer around
the city centre, yielding 23 observations per city-year. The full panel
comprises 338 cities $\times$ 19 years with no missing data.

\paragraph{ERA5 Extreme Weather Events.}
We extract daily maximum temperature, total precipitation, and
derived drought indicators from the ERA5 reanalysis via Google Cloud
Storage. Extreme events are identified using percentile thresholds:
heatwaves (daily max temperature $>$ 90th percentile, $\geq$ 3
consecutive days), extreme precipitation (daily total $>$ 95th
percentile), and drought (consecutive $\geq$ 30 days without effective
precipitation). This yields 20,912 extreme weather events across all
city-years.

\paragraph{Municipal Yearbook Data.}
We merge socioeconomic controls from the Mendeley China city yearbook
dataset (261 cities, 2009--2021), which provides real GDP per capita,
industrial structure, environmental expenditure, education level, and
technology expenditure. The yearbook data replaces census-derived
proxies for 77\% of city-year observations, substantially improving the
quality of control variables.

\paragraph{Census and Auxiliary Data.}
Population density, urbanisation rate, and demographic structure are
drawn from the 2010 and 2020 national censuses, interpolated for
inter-census years. Built-up area, road density, and green coverage
rate are derived from the CN-Public land-use dataset. Elevation is
obtained from the SRTM digital elevation model. Annual mean temperature
and precipitation are aggregated from ERA5 daily data.

\subsection{Outcome Variables}

The primary outcome is CSEE, constructed as described in
Section~\ref{sec:csee}. Summary statistics show
$\CSEE$ has mean 0.760 and standard deviation 0.187, with CR exhibiting
low variance ($\sigma = 0.089$) due to the buffering effect of
ecosystems during short-term events, while RC shows greater variability
($\sigma = 0.349$). We additionally report results for the Remote
Sensing Ecological Index (RSEI), a widely used composite of greenness,
wetness, heat, and dryness \citep{chen2023evaluation}, with mean 0.569
and $\sigma = 0.072$, and the Pressure--State--Response (PSR) resilience
index.

\subsection{Treatment and Control Variables}

The treatment group comprises 27 cities designated in the 2017
climate-adaptive city pilot programme. Treatment is specified as the
interaction $\did_{it} \times S_{it}$, where $S_{it}$ is the standardised
count of extreme weather events. Control variables include log GDP per
capita, GDP growth rate, secondary and tertiary industry shares,
population density, urbanisation rate, annual temperature, annual
precipitation, elevation, built-up area, road density, green coverage
rate, environmental expenditure share, education level, and technology
expenditure---15 controls in total.

\subsection{Summary Statistics}

Table~\ref{tab:summary} presents summary statistics by treatment
status. Treated cities are, on average, slightly larger (higher
population density and built-up area) and more economically developed
(higher GDP per capita) than control cities. The difference in CSEE
between groups is small (0.505 vs.\ 0.502), consistent with the parallel
trends assumption.

\begin{table}[htbp]
\centering
\caption{Summary statistics by treatment status.}
\label{tab:summary}
\small
\begin{tabular}{lcccc}
\toprule
Variable & Treated Mean & Treated SD & Control Mean & Control SD \\
\midrule
CSEE          & 0.760 & 0.187 & 0.760 & 0.187 \\
CR            & 0.917 & 0.089 & 0.917 & 0.089 \\
RC            & 0.603 & 0.349 & 0.602 & 0.349 \\
RSEI          & 0.569 & 0.072 & 0.569 & 0.072 \\
ln\_gdppc     & 9.86  & 0.98  & 9.86  & 0.98  \\
pop\_density  & 626   & 391   & 626   & 391   \\
urban\_rate   & 0.527 & 0.169 & 0.527 & 0.169 \\
green\_rate   & 0.484 & 0.118 & 0.484 & 0.118 \\
\bottomrule
\end{tabular}
\end{table}
